\section{习题与解答}
\begin{exercise}\label{ex:m3electricdominant}
正交均匀场$\vct E=E\hat{\vct y}$、$\vct B=B\hat{\vct z}$满足$E>cB$。带正电粒子初始实验室速度为$V\hat{\vct x}$，其中$V=c^2B/E$。求其轨迹，忽略辐射和边界。
\Diagram{fig:m3edominantq}{电占优场可变为纯电场，初速度选为相应推动速度}{\draw[axis](0,0)--(3.7,0)node[right]{$V$};\draw[axis](0,0)--(0,2.6)node[above]{$E$};\draw(1.8,1)circle(.15);\fill(1.8,1)circle(1pt);\node[right]at(2.1,1){$B$};}
\end{exercise}
\begin{solution}
图\ref{fig:m3edominantq}沿$V$推动后$B'=0$，$E'=\sqrt{E^2-c^2B^2}$，粒子初始静止。设$\alpha=qE'/m$，沿用习题\ref{ex:m1fourwork}已逐项积分的恒力解，并将加速轴改为$y'$，得到
\[
 ct'=\frac{c^2}\alpha\sinh(\alpha\tau/c),\quad
 y'=\frac{c^2}\alpha[\cosh(\alpha\tau/c)-1],\quad x'=0.
\]
逆推动给$t=\gamma t'$、$x=Vt$、$y=y'$，消去固有时得到
\[
 y(t)=\frac{c^2}\alpha\left[\sqrt{1+(\alpha t/(\gamma c))^2}-1\right].
\]
\Diagram{fig:m3edominanta}{轨迹按$X=\alpha x/(\gamma Vc)$、$Y=\alpha y/c^2$归一化，虚线为渐近线}{\draw[axis](0,0)--(4.5,0)node[right]{$X$};\draw[axis](0,0)--(0,3.5)node[above]{$Y$};\draw[curve,domain=0:4,samples=80]plot(\x,{sqrt(1+\x*\x)-1});\draw[guide](1,0)--(4.2,3.2);\node[right]at(3.7,2.4){$Y=X-1$};}
图\ref{fig:m3edominanta}的长时轨迹趋于直线斜向传播。横向速度为$v_y=(c/\gamma)\tanh(\alpha\tau/c)$，所以$v^2=V^2+c^2\tanh^2(\alpha\tau/c)/\gamma^2<c^2$。$v_x$恒定不代表$x$向动量恒定，因为总$\gamma_v$仍在增加；磁力的$x$分量正好提供这部分动量变化。
\end{solution}

\begin{exercise}\label{ex:m3decaycone}
静质量$M>2m$的粒子衰变为两个相同质量$m$的粒子，母粒子实验室速度为$V$。求子粒子在质心系的速率$u_*$；当$V>u_*$时，求所有实验室出射方向被限制的最大半角。
\Diagram{fig:m3decayconeq}{母粒子系两子粒子动量等大反向}{\draw[axis](-2.5,0)--(3.5,0)node[right]{$V$};\draw[->,thick](0,0)--(1.4,1.8)node[right]{$m$};\draw[->,thick](0,0)--(-1.4,-1.8)node[left]{$m$};}
\end{exercise}
\begin{solution}
图\ref{fig:m3decayconeq}中，母粒子静止系总动量为零，两子粒子能量相同，故$E_*=Mc^2/2$。质壳关系给
\[
 p_*c=\sqrt{M^2c^4/4-m^2c^4},\qquad
 u_*=\frac{p_*c^2}{E_*}=c\sqrt{1-4m^2/M^2}.
\]
令$z=\cos\theta_*$，母粒子沿$x$以$V$运动。速度变换的共同分母为$1+Vu_*z/c^2>0$，相除后
\[
 \tan\theta=\frac{u_*\sqrt{1-z^2}}{\gamma(V+u_*z)}.
\]
在$V>u_*$时分母始终正，所有非零速度都向前。
\Diagram{fig:m3decayconea}{亚光速子粒子可全部集中于一个前向圆锥}{\draw[axis](0,0)--(4.4,0)node[right]{$V$};\draw[curve](0,0)--(3,1.4);\draw[curve](0,0)--(3,-1.4);\draw[guide](3,-1.4)--(3,1.4);\node[above]at(2.2,1.2){$\theta_{\max}$};}
对$-1<z<1$求对数导数，最大角满足
\[
 -\frac z{1-z^2}-\frac{u_*}{V+u_*z}=0
 \quad\Longrightarrow\quad Vz+u_*=0.
\]
取$z=-u_*/V$，则$\sqrt{1-z^2}=\sqrt{V^2-u_*^2}/V$、$V+u_*z=(V^2-u_*^2)/V$，代入得到图\ref{fig:m3decayconea}的
\[
 \tan\theta_{\max}=\frac{u_*}{\gamma\sqrt{V^2-u_*^2}}.
\]
两端$z=\pm1$均有$\theta=0$，内点驻值因而为全局最大。速度相等时则出现退化：$V=u_*$时，正后方发射的子粒子在实验室静止，方向未定义；其余均向前，出射角可趋近但达不到$\pi/2$。$V<u_*$时取$z=-1$便有负的纵向速度，出现后向出射，且可有横向出射。光子$u_*=c$不可能满足$V>u_*$，因而不能用这个有限前向圆锥代替全角光行差。
\end{solution}

\begin{exercise}\label{ex:m3fizeau}
透明介质在自身静止系的非色散折射率为$n$，光沿介质运动方向传播，介质相对实验室速度$V$。求实验室相速度，并展开至$V/c$一阶。
\Diagram{fig:m3fizeauq}{介质系光速为$c/n$，介质整体沿同向运动}{\draw[curve](0,-.6)rectangle(4,1);\draw[axis](.4,.2)--(2.8,.2)node[above]{$c/n$};\draw[axis](4.2,.2)--(5.5,.2)node[above]{$V$};}
\end{exercise}
\begin{solution}
图\ref{fig:m3fizeauq}的波面速度沿$x$，由相位变换或速度合成都有$v_{\rm ph}=(c/n+V)/(1+V/(nc))$。低速把分母倒数展开为$1-V/(nc)+O(V^2/c^2)$，得$v_{\rm ph}=c/n+V(1-1/n^2)+O(V^2/c)$。
\Diagram{fig:m3fizeaua}{拖曳系数在$n=1$时为零}{\draw[axis](0,0)--(4.5,0)node[right]{$n$};\draw[axis](0,0)--(0,2.8)node[above]{$1-1/n^2$};\draw[curve,domain=1:4,samples=70]plot(\x,{2.4*(1-1/(\x*\x))});\node[below]at(1,0){$1$};}
图\ref{fig:m3fizeaua}在$n=1$时恢复真空光速$c$。本题采用非色散近似；真实色散介质中，变换也改变频率，应在相应介质静止频率处取$n$，一阶展开还会出现色散修正。
\end{solution}

\begin{exercise}\label{ex:m3currentrest}
某区域总电荷密度为零、电流密度$\vct j=j\hat{\vct x}\ne0$。是否存在使总电流为零的惯性系？解释为什么四电流不总能写为某个实固有密度乘四速度。
\Diagram{fig:m3currentq}{等量正负电荷相反运动，密度相消而电流相加}{\foreach\x in{0,1,2,3}{\node at(\x,.5){$+$};\node at(\x,-.5){$-$};}\draw[axis](.5,1)--(2,1);\draw[axis](2,-1)--(.5,-1);}
\end{exercise}
\begin{solution}
图\ref{fig:m3currentq}中，总四电流为$J^\mu=(0,j,0,0)$，不变量$J^2=-j^2<0$。假设某惯性系有$\vct j'=0$，则$J'^2=c^2\rho'^2\geq0$，与不变量符号矛盾，所以任何方向的惯性变换都不能消去总电流。

若仅沿$x$变换，可逐项看到$\rho'=-\gamma Vj/c^2$、$j'_x=\gamma j$。电中性不保持，但电流仍不为零。
\Diagram{fig:m3currenta}{类空四电流无法通过亚光速变换转到纯时间方向}{\draw[axis](-2,0)--(2.6,0)node[right]{$j_x$};\draw[axis](0,-1.8)--(0,2.4)node[above]{$c\rho$};\draw[guide](-1.8,-1.8)--(2,2);\draw[guide](-1.8,1.8)--(1.8,-1.8);\draw[->,thick](0,0)--(2,0);}
图\ref{fig:m3currenta}的四电流位于光锥外。单一带电成分共同运动时，$J^\mu=\rho_0U^\mu$且$J^2=\rho_0^2c^2\geq0$；中性电流却可能由正负两组四流相加。例如实验室密度为$+\rho$的正电荷静止，密度为$-\rho$的负电荷以$-u\hat{\vct x}$运动，总电荷为零、总电流为$+\rho u\hat{\vct x}$。各组分别有自己的静止系，总和没有共同电流静止系。

更一般地，只有$J^2>0$的非零四电流才可变到纯时间方向；$J^2<0$可变到纯空间方向，$J^2=0$且非零时两种静止化都不可行。总四电流按带符号电荷加权，与物质总能量动量不同；其类空性质不要求任何载体超光速。
\end{solution}

\begin{exercise}\label{ex:m3obliquewave}
平面波实验室传播方向与$x$成60度，观察系以$0.6c$沿$x$运动。求频率比和新方向角，核对横向波矢不变。
\Diagram{fig:m3obliquewaveq}{斜向波矢的纵横分解}{\draw[axis](0,0)--(3.8,0)node[right]{$x$};\draw[axis](0,0)--(1.5,2.598)node[right]{$\vct k$};\draw[guide](1.5,0)--(1.5,2.598);\node at(.9,.4){60度};}
\end{exercise}
\begin{solution}
图\ref{fig:m3obliquewaveq}有$\gamma=1.25$、$\cos\theta=.5$，频率比$\omega'/\omega=1.25(1-.6\times.5)=.875$。方向余弦$\cos\theta'=(.5-.6)/(1-.3)=-1/7$，所以$\theta'\simeq98.21^\circ$，略向后。
\Diagram{fig:m3obliquewavea}{两系采用相同波矢比例，横向分量的高度相同}{
\draw[axis](-1.8,0)--(2.2,0)node[right]{$k'_x$};\draw[axis](0,0)--(0,3.2)node[above]{$k'_y$};
\draw[->,thick](0,0)--(-.375,2.598)node[above left]{$\vct k'$};
\draw[guide](-.375,0)--(-.375,2.598)--(1.5,2.598)--(1.5,0);
\draw[->,dashed](0,0)--(1.5,2.598)node[above right]{$\vct k$};
\node[right]at(1.6,1.7){$k'_y=k_y$};}
新系$\sin\theta'=\sqrt{48}/7=4\sqrt3/7$，因此$k'_y=(\omega'/c)\sin\theta'=(7/8)(\omega/c)(4\sqrt3/7)=\sqrt3\omega/(2c)=k_y$。图\ref{fig:m3obliquewavea}的角度必须用负余弦判断象限，不能把正弦反求成81.79度。
\end{solution}

\begin{exercise}\label{ex:m3recoilphoton}
静止质量$M$的理想镜被能量$E$的单个光子正向撞击，光子反向弹回，镜的静质量不变。求反射光能量和镜反冲动能。
\Diagram{fig:m3recoilq}{有限质量镜不能视为预先给定速度的无限重反射面}{\draw[very thick](2,-1)--(2,1.2);\draw[->](0,.5)--(1.9,.5)node[midway,above]{$E$};\draw[->](1.9,-.3)--(0,-.3)node[midway,below]{$E_r$};\node[right]at(2.3,0){$M$};}
\end{exercise}
\begin{solution}
图\ref{fig:m3recoilq}镜末能量$Mc^2+E-E_r$，末动量$(E+E_r)/c$。代入镜的质壳关系，
\[
 (Mc^2+E-E_r)^2-(E+E_r)^2=M^2c^4.
\]
展开得$2Mc^2(E-E_r)=4EE_r$，所以$E_r=E/[1+2E/(Mc^2)]$。
\Diagram{fig:m3recoila}{镜的反冲动能等于光子损失能量}{\draw[axis](0,0)--(4.3,0)node[right]{$E/(Mc^2)$};\draw[axis](0,0)--(0,2.6)node[above]{$E_r/E$};\draw[curve,domain=0:3.5,samples=70]plot(\x,{2.2/(1+2*\x)});}
反冲动能$K_M=E-E_r=2E^2/(Mc^2+2E)$。图\ref{fig:m3recoila}在$M\to\infty$趋无频移；有限质量时必红移。这里不能先选某个“镜速”代入两次多普勒式，因为速度在单次碰撞中改变，需用全程守恒确定。
\end{solution}

\begin{exercise}\label{ex:m3invariantfields}
在某系有$\vct E=E_0\hat{\vct y}$、$\vct B=2E_0\hat{\vct z}/c$。求消电场的速度、剩余磁场和该系中粒子静止时对应实验室的速度。
\Diagram{fig:m3invariantq}{磁占优正交场}{\draw[axis](0,0)--(3.5,0)node[right]{$x$};\draw[axis](0,0)--(0,2.1)node[above]{$E_0$};\draw(1.6,1)circle(.15);\fill(1.6,1)circle(1pt);\node[right]at(1.9,1){$2E_0/c$};}
\end{exercise}
\begin{solution}
图\ref{fig:m3invariantq}给出$E=E_0$、$B=2E_0/c$，且$\vct E\perp\vct B$，磁占优。沿$+x$变换，有$E_y'=\gamma(E_0-V\,2E_0/c)$。令其为零，得到$V=c/2$、$\gamma=2/\sqrt3$。再求剩余磁场，
\[
 B_z'=\frac2{\sqrt3}\left(\frac{2E_0}c-\frac{E_0}{2c}\right)
 =\frac{\sqrt3E_0}c.
\]
也可由$B'^2=B^2-E^2/c^2=3E_0^2/c^2$独立核对其大小，方向由具体变换保留为$+z$。
\Diagram{fig:m3invarianta}{消场系中静止对应实验室匀速漂移}{\fill(0,0)circle(2pt);\draw[axis](0,.5)--(2.5,.5)node[right]{$V=c/2$};\node[below]at(1,-.4){$\vct E+\vct V\times\vct B=0$};}
图\ref{fig:m3invarianta}中的粒子在新系静止，逆速度变换给实验室$\vct v=(c/2)\hat{\vct x}$。其磁力方向为$\hat{\vct x}\times\hat{\vct z}=-\hat{\vct y}$，大小$q(c/2)(2E_0/c)=qE_0$，正好与电力相抵。

若只要求新系粒子沿$z'$运动而非静止，则逆变换仍有$v_x=c/2$，还允许$v_z=v_z'/\gamma$，范围$|v_z|<\sqrt3c/2$。这进一步说明消电场不等于消去所有粒子的力：沿$\vct B'$运动仍零力，垂直$\vct B'$运动却受磁力并回旋。参考系选择由场决定，具体粒子是否无偏转还要代入其速度。
\end{solution}

\begin{exercise}\label{ex:m3properpower}
粒子静质量不变，某实验室三力始终平行瞬时速度。证明四力模等于三力模，并说明这并不意味着四力时间分量为零。
\Diagram{fig:m3properq}{三力与速度平行而功率非零}{\draw[axis](0,0)--(3.5,0)node[right]{$\vct v$};\draw[axis](0,.6)--(2.5,.6)node[right]{$\vct f$};}
\end{exercise}
\begin{solution}
图\ref{fig:m3properq}中，取瞬时速度沿$x$正向，平行三力的有向分量为$f$。四力为
\[
 F^\mu=\gamma_v(\beta_v f,f,0,0),\qquad
 F^2=\gamma_v^2f^2(\beta_v^2-1)=-f^2.
\]
故固有三力大小$f_0=\sqrt{-F^2}=|f|$。这里所谓四力模是类空四力的$\sqrt{-F^2}$，不是四个分量平方和的欧氏模。
\Diagram{fig:m3propera}{类空四力可有非零时间分量；零功率参考系中的时间分量为零}{\draw[axis](0,0)--(3.5,0)node[right]{$F^x$};\draw[axis](0,0)--(0,2.6)node[above]{$F^0$};\draw[->,thick](0,0)--(2.8,1.4);\draw[guide](0,0)--(2.4,2.4);\node[right]at(2.8,1.4){$F^0/F^x=v/c$};}
图\ref{fig:m3propera}中$F^0=\gamma_vfv/c$，只要$fv\ne0$便非零。变到粒子瞬时静止系后，
\[
 F_*^0=\gamma_v(F^0-\beta_vF^x)=0,\qquad F_*^x=f.
\]
但零时间分量并非瞬时静止系独有。以该静止系为起点，再沿$y$以任意$w<c$变换：原四力$(0,f,0,0)$的$0,y$分量均为零，变换后仍是$(0,f,0,0)$；粒子却以$-w\hat{\vct y}$运动。此系三力为$(f/\gamma_w,0,0)$，与粒子速度垂直，功率仍为零。

一般判据是$F^0=\gamma\vct f\cdot\vct v/c=0$，即该系中瞬时功率为零。四维正交$u\cdot F=0$则在所有参考系成立，它不要求每个参考系中的时间分量都为零。
\end{solution}

\begin{exercise}\label{ex:m3threshold}
静质量$m$的入射粒子撞击同样静质量的静止靶粒子，末态需产生四个质量均为$m$的粒子。忽略其他限制，求入射动能阈值，并与两束等能反向对撞比较。
\Diagram{fig:m3thresholdq}{固定靶与对撞的入射配置}{\draw[->](0,.5)--(2,.5);\fill(2.8,.5)circle(2pt);\node[above]at(1,.5){$E$};\node[right]at(3,.5){固定靶};\draw[->](0,-1)--(1.7,-1);\draw[->](3.6,-1)--(1.9,-1);\node[right]at(4,-1){对撞};}
\end{exercise}
\begin{solution}
在图\ref{fig:m3thresholdq}的固定靶实验室系，入射总能量记为$E$，动量记为$p$，有$E^2-p^2c^2=m^2c^4$。初态总能量$E+mc^2$、总动量$p$，因此
\begin{align*}
 s&=(E+mc^2)^2-p^2c^2\\
 &=E^2-p^2c^2+2mc^2E+m^2c^4\\
 &=2m^2c^4+2mc^2E.
\end{align*}
阈值时四个末态粒子在质心系相对静止，最小总能量为$4mc^2$，故$s=16m^2c^4$，解得$E=7mc^2$、$K=6mc^2$。
\Diagram{fig:m3thresholda}{阈值末态在质心系中无相对运动}{\foreach\x in{0,1,2,3}{\fill(\x,0)circle(3pt)node[below]{$m$};}\draw[curve](-.4,-.5)rectangle(3.4,.6);\node[above]at(1.5,.6){总静能$4mc^2$};}
图\ref{fig:m3thresholda}的末态静止是指质心系。实验室的初态总能量为$8mc^2$，总动量为$\sqrt{7^2-1}\,mc=4\sqrt3\,mc$，因而质心速度$V=Pc^2/E_T=\sqrt3c/2$、$\gamma_V=2$。末态整体动能为
\[
 K_{\rm bulk}=(\gamma_V-1)\,4mc^2=4mc^2.
\]
六份入射动能中，$2mc^2$用来增加末态静能，$4mc^2$用于整体运动。若改成两束等能反向对撞，总动量一开始为零，阈值只需每束总能量$2mc^2$，每束动能$mc^2$，两束共$2mc^2$。两种装置的末态最小静能相同，固定靶所需额外能量来自总动量约束。
\end{solution}

\begin{exercise}\label{ex:m3compton}
能量$E$的光子撞击静止电子，散射光子与入射方向成$\theta$角。用四动量守恒推导波长改变量。
\Diagram{fig:m3comptonq}{散射光子与反冲电子}{\draw[->,thick](-2,0)--(0,0)node[midway,above]{$E$};\draw[->,thick](0,0)--(2,1.7)node[right]{$E'$};\draw[->](0,0)--(2,-1.1)node[right]{电子};\draw[guide](0,0)--(2.5,0);\node at(.8,.35){$\theta$};}
\end{exercise}
\begin{solution}
图\ref{fig:m3comptonq}中，初始电子四动量为$P=(mc,\vct0)$，初末光子四动量记为$k,k'$，其能量分别为$E,E'$。末电子$P'=P+k-k'$，电子静质量不变，故
\[
 0=P'^2-P^2=2P\cdot k-2P\cdot k'-2k\cdot k'.
\]
光子各自平方为零，两个光子的空间夹角为$\theta$，因此
\[
 P\cdot k=mE,\quad P\cdot k'=mE',\quad
 k\cdot k'=\frac{EE'}{c^2}(1-\cos\theta).
\]
代入并除以$EE'$，得$1/E'-1/E=(1-\cos\theta)/(mc^2)$。
\Diagram{fig:m3comptona}{散射角越大，光子能量损失越大}{\draw[axis](0,0)--(4.1,0)node[right]{$\theta$};\draw[axis](0,0)--(0,2.8)node[above]{$\lambda'-\lambda$};\draw[curve,domain=0:180,samples=70]plot({\x/50},{1-cos(\x)});\node[below]at(3.6,0){$\pi$};\node[left]at(0,2){$2h/(mc)$};}
以$E=hc/\lambda$表示，即为图\ref{fig:m3comptona}的康普顿关系
\[
 \lambda'-\lambda=\frac h{mc}(1-\cos\theta),\qquad
 E'=\frac E{1+(E/mc^2)(1-\cos\theta)}.
\]
电子反冲不能在守恒中省掉。取散射光子在$x,y$平面且$y$分量为正，末电子分量为
\[
 p_{ex}=\frac{E-E'\cos\theta}c,\qquad
 p_{ey}=-\frac{E'\sin\theta}c,
 \quad K_e=E-E'.
\]
这些分量应与题图中的反冲方向一致。由它们平方得到$p_e^2c^2=E^2+E'^2-2EE'\cos\theta$；结合$K_e$又有$p_e^2c^2=K_e^2+2mc^2K_e$，两式相等恰恢复上面的散射能量式。正向散射无能量转移，反向散射波长改变量最大为$2h/(mc)$。
\end{solution}

\begin{exercise}\label{ex:m3circular}
相对论带电粒子在均匀磁场中作圆周运动，速度垂直磁场。忽略辐射。求半径、实验室角频率、固有加速度及瞬时静止系电场。
\Diagram{fig:m3circleq}{垂直磁场中的圆周运动，取正电荷}{\draw[curve](0,0)circle(1.5);\fill(1.5,0)circle(2pt);\draw[axis](1.5,0)--(1.5,1.2)node[right]{$v$};\draw[axis](1.5,0)--(.3,0)node[above]{$qvB$};\draw(0,0)circle(.15);\draw(-.1,-.1)--(.1,.1);\draw(-.1,.1)--(.1,-.1);\node[below]at(0,-.3){$B$入纸面};}
\end{exercise}
\begin{solution}
磁力始终垂直速度，$\dd E/\dd t=0$，故$\gamma_v$恒定。图\ref{fig:m3circleq}中动量转向率$p v/R=|q|vB$，因此$R=p/(|q|B)=\gamma_vmv/(|q|B)$，$\omega=v/R=|q|B/(\gamma_vm)$。
\Diagram{fig:m3circlea}{局部瞬时静止系中的电场指向圆心}{\fill(0,0)circle(2pt)node[below]{粒子};\draw[axis](0,0)--(-2,0)node[above]{$\vct E_0=\gamma_v\vct v\times\vct B$};\draw[guide](0,-1)--(0,1.2);\node[right]at(.2,.7){$\vct v$切向};}
实验室向心加速度$a=v^2/R$属垂直分量，所以固有加速度$\alpha=\gamma_v^2v^2/R=\gamma_v|q|vB/m$。瞬时静止系场变换给图\ref{fig:m3circlea}的$\vct E_0=\gamma_v\vct v\times\vct B$；静止粒子所受电力大小$|q|\gamma_vvB=m\alpha$，与加速度关系相符。该系电场可做瞬时加速，却不意味着实验室磁场做功。
\end{solution}

\begin{exercise}\label{ex:m3sheet}
无限带电平面静止系面密度为$\sigma_0$，平面为$xy$面。观察系沿$x$以$V$运动。求平面两侧的电磁场，并用变换后的面电流核对磁场方向。
\Diagram{fig:m3sheetq}{静止带电平面两侧的电场}{\draw[curve](-2,0)--(2,0);\foreach\x in{-1.5,-.5,.5,1.5}{\node at(\x,0){$+$};\draw[axis](\x,.2)--(\x,1.2);\draw[axis](\x,-.2)--(\x,-1.2);}\node[right]at(2,0){$xy$面};\node[right]at(1.5,1.2){$z>0$};}
\end{exercise}
\begin{solution}
图\ref{fig:m3sheetq}中$E_z=\operatorname{sgn}(z)\sigma_0/(2\varepsilon_0)$、$B=0$。沿$x$变换得$E_z'=\gamma E_z$、$B_y'=\gamma VE_z/c^2$。面密度$\sigma'=\gamma\sigma_0$，平面在新系沿$-x$运动，面电流$K_x'=-\gamma\sigma_0V$。
\Diagram{fig:m3sheeta}{新系面电流沿$-x$，上侧磁场沿$+y$}{\draw[axis](-2,0)--(2,0)node[right]{$x'$};\draw[axis](0,-1.7)--(0,1.8)node[above]{$z'$};\draw[very thick,->](1,0)--(-1,0)node[below]{$\vct K'$};\draw(.8,1)circle(.15);\draw(.7,.9)--(.9,1.1);\draw(.7,1.1)--(.9,.9);\node[right]at(1.1,1){$B_y'$入纸面};}
按图\ref{fig:m3sheeta}的右手定则，无限面电流上侧$B_y'=-\mu_0K_x'/2=\mu_0\gamma\sigma_0V/2$，与$\gamma VE_z/c^2$完全一致，因为$\varepsilon_0\mu_0c^2=1$。上下侧场都反号。这里$x'$向右、$z'$向上时，$+y'$朝纸内，基底方向需在图中明确。
\end{solution}

\begin{exercise}\label{ex:m3gauge}
真空平面波的四势为$A^\mu=\epsilon^\mu\ee^{-ik_\nu x^\nu}$，满足洛伦兹规范。求极化四矢量的约束，并说明沿波矢的四势分量为何不改变场。
\Diagram{fig:m3gaugeq}{横向极化与传播方向}{\draw[axis](0,0)--(4,0)node[right]{$\vct k$};\draw[axis](1,0)--(1,1.8)node[above]{$\vct\epsilon_\perp$};\draw[guide](0,1)--(4,1);\node[below]at(2,-.4){等相位波面垂直$\vct k$};}
\end{exercise}
\begin{solution}
规范条件$\partial_\mu A^\mu=0$给$k_\mu\epsilon^\mu=0$。若再选无标势的横向规范$\epsilon^0=0$，则$\vct k\cdot\vct\epsilon=0$，得到图\ref{fig:m3gaugeq}的两个独立横向极化方向。
\Diagram{fig:m3gaugea}{规范变化沿四维波矢方向，横向物理场不变}{\draw[axis](0,0)--(4.3,0)node[right]{$k^\mu$方向};\draw[axis](0,0)--(0,2.2)node[above]{横向};\draw[->,thick](0,0)--(1.5,1.6)node[above]{$\epsilon^\mu$};\draw[->,thick](0,0)--(3.2,1.6)node[right]{$\epsilon^\mu+ak^\mu$};\draw[guide](1.5,1.6)--(3.2,1.6);}
场张量由反对称组合$k^\mu\epsilon^\nu-k^\nu\epsilon^\mu$给出。把图\ref{fig:m3gaugea}中$\epsilon^\mu$改为$\epsilon^\mu+ak^\mu$，新增项$a(k^\mu k^\nu-k^\nu k^\mu)=0$，所以场不变。并且$k^2=0$使$k\cdot\epsilon$约束仍成立。图为四维分量关系示意，不是普通空间中把电场沿传播方向倾斜。
\end{solution}

\begin{exercise}\label{ex:m3parallelfields}
某处电磁场平行且均非零。能否通过惯性变换同时消去其中一个场？若两场本来不平行，两个场不变量可以提供哪些必要限制？
\Diagram{fig:m3parallelq}{平行电磁场的点积不为零}{\draw[axis](0,0)--(0,1.65)node[above]{$\vct E$};\draw[axis](1.2,0)--(1.2,1.35)node[above]{$\vct B$};\node[below]at(.6,0){同向或反向均可};}
\end{exercise}
\begin{solution}
图\ref{fig:m3parallelq}中$\vct E\cdot\vct B=\pm EB\ne0$，它是洛伦兹不变量。若在某系消去任意一个场，点积就为零，矛盾。因此不存在这种惯性系。
\Diagram{fig:m3parallela}{消场必须先满足两个不变量的限制}{\draw[axis](0,0)--(5,0)node[right]{$E^2-c^2B^2$};\draw[curve](2.5,-.6)--(2.5,.8);\node[below]at(2.5,-.6){$0$};\node[above]at(1,.2){纯磁场可能};\node[above]at(4,.2){纯电场可能};\node[below]at(2.5,-1.1){还须$\vct E\cdot\vct B=0$};}
一般场若可变为纯电场，必要有$\vct E\cdot\vct B=0$且$E^2-c^2B^2>0$；若可变为纯磁场，则后一不变量小于0。图\ref{fig:m3parallela}的零边界需要另行处理：两个不变量都零的非零场不能被消去，因为若某系变成零场，逆变换会使原系也为零。对正交场，例\ref{ex:m3crossed}已给出满足这些条件的具体速度。
\end{solution}
